% Généralisation de Swendsen--Wang aux motifs : sur chaque motif b, un SEUL
% tirage corrélé choisit le sous-lien gelé omega_b, sous-ensemble de b.
% Convention : plein = attraction, pointillés = répulsion, gras = satisfaite,
% gras rouge = gelée.
    \begin{tikzpicture}[x=1cm,y=1cm,scale=0.71,every node/.style={scale=0.71}]

      % ================= triangle attractif : tout ou rien =================
      \node[font=\footnotesize,align=center] at (1.85,2.20) {Triangle \textbf{attractif}\\[-1pt] {\scriptsize tout ou rien}};

      \begin{scope}[shift={(0.55,0)}]
        \node[spin] (a1) at (0,0) {$+$};
        \node[spin] (b1) at (1.3,0) {$+$};
        \node[spin] (c1) at (0.65,1.13) {$+$};
        \draw[posedge,freeze] (a1)--(b1);
        \draw[posedge,freeze] (b1)--(c1);
        \draw[posedge,freeze] (c1)--(a1);
        \node[font=\scriptsize] at (0.65,-0.55) {$\omega_b=\varnothing$};
      \end{scope}

      \begin{scope}[shift={(2.75,0)}]
        \node[spin] (a2) at (0,0) {$+$};
        \node[spin] (b2) at (1.3,0) {$+$};
        \node[spin] (c2) at (0.65,1.13) {$+$};
        \draw[posedge,freeze,inria-rouge] (a2)--(b2);
        \draw[posedge,freeze,inria-rouge] (b2)--(c2);
        \draw[posedge,freeze,inria-rouge] (c2)--(a2);
        \node[font=\scriptsize] at (0.65,-0.55) {$\omega_b=b$};
      \end{scope}

      % ============ triangle répulsif frustré : jamais les trois arêtes ====
      \node[font=\footnotesize,align=center] at (7.5,2.20) {Triangle \textbf{répulsif} (frustré)\\[-1pt] {\scriptsize jamais les trois arêtes}};

      \begin{scope}[shift={(5.6,0)}]
        \node[spin] (a3) at (0,0) {$+$};
        \node[spin] (b3) at (1.3,0) {$+$};
        \node[spin] (c3) at (0.65,1.13) {$-$};
        \draw[negedge,line width=0.6pt] (a3)--(b3);
        \draw[negedge,freeze] (b3)--(c3);
        \draw[negedge,freeze] (c3)--(a3);
        \node[font=\scriptsize] at (0.65,-0.55) {$\omega_b=\varnothing$};
      \end{scope}

      \begin{scope}[shift={(7.8,0)}]
        \node[spin] (a4) at (0,0) {$+$};
        \node[spin] (b4) at (1.3,0) {$+$};
        \node[spin] (c4) at (0.65,1.13) {$-$};
        \draw[negedge,line width=0.6pt] (a4)--(b4);
        \draw[negedge,freeze,inria-rouge] (b4)--(c4);
        \draw[negedge,freeze] (c4)--(a4);
        \node[font=\scriptsize] at (0.65,-0.55) {$|\omega_b|=1$};
      \end{scope}

    \end{tikzpicture}
